\documentclass[10pt,a4paper]{amsart}
\usepackage{amsmath,amssymb,amsthm,graphicx}
\usepackage{verbatim}
\newenvironment{reference}{\comment}{\endcomment}
%\newenvironment{reference}{}{}
\newenvironment{slogan}{\comment}{\endcomment}
\newenvironment{history}{\comment}{\endcomment}

% For commutative diagrams we use Xy-pic
\usepackage[all]{xy}

% We use 2cell for 2-commutative diagrams.
\xyoption{2cell}
% Theorem environments.
%
\theoremstyle{plain}
\newtheorem{theorem}[subsection]{Theorem}
\newtheorem{proposition}[subsection]{Proposition}
\newtheorem{lemma}[subsection]{Lemma}

\theoremstyle{definition}
\newtheorem{definition}[subsection]{Definition}
\newtheorem{example}[subsection]{Example}
\newtheorem{exercise}[subsection]{Exercise}
\newtheorem{situation}[subsection]{Situation}

\theoremstyle{remark}
\newtheorem{remark}[subsection]{Remark}
\newtheorem{remarks}[subsection]{Remarks}

\numberwithin{equation}{subsection}

% Macros
%
\def\lim{\mathop{\mathrm{lim}}\nolimits}
\def\colim{\mathop{\mathrm{colim}}\nolimits}
\def\Spec{\mathop{\mathrm{Spec}}}
\def\Hom{\mathop{\mathrm{Hom}}\nolimits}
\def\Ext{\mathop{\mathrm{Ext}}\nolimits}
\def\SheafHom{\mathop{\mathcal{H}\!\mathit{om}}\nolimits}
\def\SheafExt{\mathop{\mathcal{E}\!\mathit{xt}}\nolimits}
\def\Sch{\mathit{Sch}}
\def\Mor{\mathop{\mathrm{Mor}}\nolimits}
\def\Ob{\mathop{\mathrm{Ob}}\nolimits}
\def\Sh{\mathop{\mathit{Sh}}\nolimits}
\def\NL{\mathop{N\!L}\nolimits}
\def\CH{\mathop{\mathrm{CH}}\nolimits}
\def\proetale{{pro\text{-}\acute{e}tale}}
\def\etale{{\acute{e}tale}}
\def\QCoh{\mathit{QCoh}}
\def\Ker{\mathop{\mathrm{Ker}}}
\def\Im{\mathop{\mathrm{Im}}}
\def\Coker{\mathop{\mathrm{Coker}}}
\def\Coim{\mathop{\mathrm{Coim}}}

% Boxtimes
%
\DeclareMathSymbol{\boxtimes}{\mathbin}{AMSa}{"02}

%
% Macros for moduli stacks/spaces
%
\def\QCohstack{\mathcal{QC}\!\mathit{oh}}
\def\Cohstack{\mathcal{C}\!\mathit{oh}}
\def\Spacesstack{\mathcal{S}\!\mathit{paces}}
\def\Quotfunctor{\mathrm{Quot}}
\def\Hilbfunctor{\mathrm{Hilb}}
\def\Curvesstack{\mathcal{C}\!\mathit{urves}}
\def\Polarizedstack{\mathcal{P}\!\mathit{olarized}}
\def\Complexesstack{\mathcal{C}\!\mathit{omplexes}}
% \Pic is the operator that assigns to X its picard group, usage \Pic(X)
% \Picardstack_{X/B} denotes the Picard stack of X over B
% \Picardfunctor_{X/B} denotes the Picard functor of X over B
\def\Pic{\mathop{\mathrm{Pic}}\nolimits}
\def\Picardstack{\mathcal{P}\!\mathit{ic}}
\def\Picardfunctor{\mathrm{Pic}}
\def\Deformationcategory{\mathcal{D}\!\mathit{ef}}
\usepackage{cleveref}
\crefname{originaltheorem}{original theorem}{original theorems}
\Crefname{originaltheorem}{Original theorem}{Original theorems}
\makeatletter
\newlabel{theorem-baer-grothendieck}{{05NX}{source}{}{sourcepointer-1}{}}
\newlabel{theorem-baer-grothendieck@cref}{{[originaltheorem][0][]05NX}{[1][0][]source}{}{}{}}
\makeatother
\crefname{originalproposition}{original proposition}{original propositions}
\Crefname{originalproposition}{Original proposition}{Original propositions}
\makeatletter
\newlabel{sets-proposition-exist-ordinals-large-cofinality}{{05N3}{source}{}{sourcepointer-2}{}}
\newlabel{sets-proposition-exist-ordinals-large-cofinality@cref}{{[originalproposition][0][]05N3}{[1][0][]source}{}{}{}}
\makeatother
\crefname{originalproposition}{original proposition}{original propositions}
\Crefname{originalproposition}{Original proposition}{Original propositions}
\makeatletter
\newlabel{proposition-objects-are-small}{{079F}{source}{}{sourcepointer-3}{}}
\newlabel{proposition-objects-are-small@cref}{{[originalproposition][0][]079F}{[1][0][]source}{}{}{}}
\makeatother
\crefname{originallemma}{original lemma}{original lemmas}
\Crefname{originallemma}{Original lemma}{Original lemmas}
\makeatletter
\newlabel{lemma-characterize-injective}{{079G}{source}{}{sourcepointer-4}{}}
\newlabel{lemma-characterize-injective@cref}{{[originallemma][0][]079G}{[1][0][]source}{}{}{}}
\makeatother
\title{A Grothendieck abelian category has functorial injective embeddings}
\author{The Stacks Project Authors}
\date{stacks-project:injectives.tex:local-snapshot-20261012; GNU FDL 1.2 or later}
\begin{document}
\maketitle
\noindent\textit{Editorial note.} All delivered mathematical text is the author's own, copied from the pinned Stacks Project chapter file sites-cohomology.tex; no statement or proof step was rewritten, reordered or removed.\par\medskip
\noindent\textit{Editorial note.} The proof is the author's complete argument; its opening sentence 'Please compare with the proof of Theorem theorem-baer-grothendieck.' is a reading aid pointing at a companion construction, not a substitute for one, since the construction and its verification are then written out in the counted span. The sentence is preserved verbatim.\par\medskip
\noindent\textit{Editorial note.} Every labelled reference occurring inside the counted spans is carried in assembly.external\_references as a non-load-bearing navigation pointer (same-chapter ones at their definition site, cross-chapter ones at the invocation inside the counted span), so no load-bearing mathematics is replaced by a stub. No bibliography citation occurs inside the counted spans.\par\medskip
\noindent\textit{Editorial note.} The same-name check for this card scanned every staged Stacks chapter of this repository -- all sources/stacks-*/pinned-original-source.tex files, the card's own chapter excluded -- and the labels found are recorded in dedup.same\_name\_cross\_chapter on the card.\par\medskip
\noindent\textit{Editorial note.} The preamble is the chapter's own retained minimal preamble (sources/stacks-sites-cohomology/retained-preamble.tex), a verbatim excerpt of the native preamble: the Xy-pic setup and the native theorem declarations with the mathematical macros.\par\medskip
\noindent\textit{Allowed standard prerequisites (editorial).}\begin{itemize}
\item The transfinite construction of an injective embedding. Same chapter, Theorem theorem-baer-grothendieck and the construction of the objects attached to an ordinal used there. Supplies the template of the transfinite construction, which the author adapts by hand and which the opening sentence of the proof invites the reader to compare with.
\item A criterion for injectivity. Same chapter, Lemma lemma-characterize-injective and Lemma lemma-characterize-injectives. Supply the criteria for injectivity of an object, and of a sheaf of abelian groups, that are applied at the end of the construction.
\item Smallness of the objects in a Grothendieck abelian category. Same chapter, the smallness results established earlier in the chapter. Supply that the transfinite process terminates, which makes the construction functorial.
\end{itemize}\medskip
\ifdefined\hypertarget\hypertarget{sourcepointer-1}{}\fi
\noindent\textit{Source reference.} Original theorem 05NX: sources/stacks-injectives/pinned-original-source.tex, source lines 567--567. Navigation pointer invoked at proof line 1723; the counted span contains the pointer only and the cited result is not reproduced inside it, so the pointer is registered as a non-load-bearing navigation stub in the established form for Stacks material. No bibliography citation is involved.\par\medskip
\ifdefined\hypertarget\hypertarget{sourcepointer-2}{}\fi
\noindent\textit{Source reference.} Original proposition 05N3: sources/stacks-sets/pinned-original-source.tex, source lines 434--434. Navigation pointer invoked at proof line 1748; the counted span contains the pointer only and the cited result is not reproduced inside it, so the pointer is registered as a non-load-bearing navigation stub in the established form for Stacks material. No bibliography citation is involved.\par\medskip
\ifdefined\hypertarget\hypertarget{sourcepointer-3}{}\fi
\noindent\textit{Source reference.} Original proposition 079F: sources/stacks-injectives/pinned-original-source.tex, source lines 1634--1634. Navigation pointer invoked at proof line 1754; the counted span contains the pointer only and the cited result is not reproduced inside it, so the pointer is registered as a non-load-bearing navigation stub in the established form for Stacks material. No bibliography citation is involved.\par\medskip
\ifdefined\hypertarget\hypertarget{sourcepointer-4}{}\fi
\noindent\textit{Source reference.} Original lemma 079G: sources/stacks-injectives/pinned-original-source.tex, source lines 1671--1671. Navigation pointer invoked at proof line 1758; the counted span contains the pointer only and the cited result is not reproduced inside it, so the pointer is registered as a non-load-bearing navigation stub in the established form for Stacks material. No bibliography citation is involved.\par\medskip
\begin{theorem}
\label{theorem-injective-embedding-grothendieck}
Let $\mathcal{A}$ be a Grothendieck abelian category.
Then $\mathcal{A}$ has functorial injective embeddings.
\end{theorem}

\begin{proof}
Please compare with the proof of
Theorem \ref{theorem-baer-grothendieck}.
Choose a generator $U$ of $\mathcal{A}$. For an object $M$ we define
$\mathbf{M}(M)$ by the following pushout diagram
$$
\xymatrix{
\bigoplus_{N \subset U}
\bigoplus_{\varphi \in \Hom(N, M)}
N \ar[r] \ar[d] & M \ar[d] \\
\bigoplus_{N \subset U}
\bigoplus_{\varphi \in \Hom(N, M)}
U \ar[r] &  \mathbf{M}(M).
}
$$
Note that $M \mapsto \mathbf{M}(M)$ is a functor and that there
exist functorial injective maps $M \to \mathbf{M}(M)$. By transfinite
induction we define functors $\mathbf{M}_\alpha(M)$ for every
ordinal $\alpha$. Namely, set $\mathbf{M}_0(M) = M$. Given
$\mathbf{M}_\alpha(M)$ set
$\mathbf{M}_{\alpha + 1}(M) = \mathbf{M}(\mathbf{M}_\alpha(M))$.
For a limit ordinal $\beta$ set
$$
\mathbf{M}_\beta(M) = \colim_{\alpha < \beta} \mathbf{M}_\alpha(M).
$$
Finally, pick any ordinal $\alpha$ whose cofinality is greater than $|U|$.
Such an ordinal exists by
Sets, Proposition \ref{sets-proposition-exist-ordinals-large-cofinality}.
We claim that $M \to \mathbf{M}_\alpha(M)$ is the desired functorial
injective embedding. Namely, if $N \subset U$ is a subobject and
$\varphi : N \to \mathbf{M}_\alpha(M)$ is a morphism, then we see that
$\varphi$ factors through $\mathbf{M}_{\alpha'}(M)$ for some
$\alpha' < \alpha$ by
Proposition \ref{proposition-objects-are-small}.
By construction of $\mathbf{M}(-)$ we see that $\varphi$ extends to
a morphism from $U$ into $\mathbf{M}_{\alpha' + 1}(M)$ and hence into
$\mathbf{M}_\alpha(M)$. By
Lemma \ref{lemma-characterize-injective}
we conclude that $\mathbf{M}_\alpha(M)$ is injective.
\end{proof}

\end{document}
